\title{Nonlinear Resolvent Geometry and Error Bounds\\
for Generalized Proximal Point Convergence}
\author{菠萝之神}
\date{}
\maketitle

\begin{abstract}
广义近端点算法以核增量代替物理位移，但核坐标的收敛本身不足以保证迭代点收敛。
本文给出一组直接在物理空间中使用的局部充分条件：
相对于输入最近零点的核反射估计、实际输出处的真实残差误差界，
以及相邻允许转移的物理位移控制。
当它们给出的距离递推严格收缩且物理步长可和时，
输入覆盖和严格初始长度预算保证全部允许选择均可无限延拓，
具有有限长度并收敛到闭零集中的一点。
恒等核准确恢复普通 PPA 的非线性反射几何与误差界定理；
保留核能量项则给出另一条收缩判据，其单调特例允许任意正步长。
两个非孤立零集模型分别展示平方根阶切向漂移的可和性，
以及四次 Bregman 核对二次目标产生的三次物理收敛。
最后，对具有全尺度次线性反射界的完整普通近端映射，
我们证明局部吸引流形的指数为其欧拉示性数，剩余零点的总指数为 $1-\chi(M)$，其中 $M$ 为该局部流形。
一个弱凸径向目标具体展示吸引极小点圆与补偿临界点的共存。
\end{abstract}

\medskip
\noindent\textbf{关键词：}广义近端点算法；Bregman 近端；
非线性反射几何；残差误差界；有限长度；非孤立零点；固定点指数。

\section{引言}
\label{sec:introduction}

给定 $F:\mathbb R^n\rightrightarrows\mathbb R^n$ 与 $\lambda>0$，
广义近端点算法通过
\begin{equation}
 v(x^k)-v(x^{k+1})\in\lambda F(x^{k+1}),\qquad
 J^v_{\lambda F}:=(v+\lambda F)^{-1}v
 \label{eq:ppa-intro}
\end{equation}
寻找 $F$ 的零点。恒等核 $v=I$ 给出普通近端点算法；
固定正定线性核给出度量近端步；梯度核 $v=\nabla h$ 则产生 Bregman 近端步。
这种核化近端关系已有系统研究
\cite{BauschkeBorweinCombettes2003,BuiCombettes2020Warped}。
本文关心一个局部问题：当反射近端几何只能以非线性模量控制，
而核与物理坐标之间也没有 Lipschitz 逆界时，
什么条件能保证全部允许迭代轨道收敛到非孤立零集中的一点？

普通 PPA 的经典分析以单调性为起点。
极大单调性保证近端映射处处单值且牢固非扩张
\cite{Rockafellar1976PPA}；度量次正则性进一步给出局部线性率
\cite[Theorem~3.1]{Leventhal2009MetricSubregularityPPA}。
非单调理论则放宽近端几何并保留适当正则性：
负余单调算子允许非精确、松弛近端方法
\cite{CombettesPennanen2004Cohypomonotone}，
近乎平均映射框架同时处理扩张与集值行为
\cite{LukeThaoTam2018Quantitative}。
Luke--Tam 将次单调性、度量次正则性与截断值域覆盖结合，
建立适当局部近端选择的线性收敛
\cite[Assumption~2, Theorem~2]{LukeTam2025GeneralizedMonotonicity}。
半单调性及 oblique weak Minty 条件还允许松弛与预条件近端迭代
\cite{EvensPasLatafatPatrinos2025SemimonotonePPA}。
本文在线性阶给出与其中二次型条件的精确参数对应，
并以输出真残差保留一步能量，恢复经典单调 PPA 的全部正步长范围。

非线性误差界为较弱的近端几何提供了另一种补偿。
H\"older 次正则性、一般模量与高阶次正则性已经发展出丰富的判据
\cite{LiMordukhovich2012HolderSubregularityPPA,
Kruger2016NonlinearSubregularity,MordukhovichOuyang2015HigherOrderSubregularity}。
对于实际广义近端输出 $y$，包含关系只给出
\[
 r_F(y):=d(0,F(y))\leq\|v(x)-v(y)\|/\lambda.
\]
误差界据此控制 $d(y,S)$，核几何则控制右侧的增量。
要由距离衰减得到点收敛，还须把核增量转为可和的物理位移。
这项控制既可来自核逆的连续模，也可由算法对核纤维内运动的限制直接给出。
因此我们保留闭物理零集 $S\subset F^{-1}(0)$，分别验证核几何、真残差与物理位移。

证明先使用一个常见的距离--长度判定：距离几何衰减与可和步长共同保证点收敛。
主定理将这两项估计从相同的允许近端关系中导出，
再用严格长度预算逐个证明有限轨道前缀留域，使覆盖可以继续使用。
允许关系可以严格多值；同一组常数控制全部允许选择，极限点可以随选择改变。
恒等核下，前向锚界与物理位移界自动成立，因而准确恢复普通 PPA 的原有条件。
全对 RL 图像条件是验证最近零点锚的一种充分方式。

两个简短模型解释上述条件的不同作用。
三角系统的法向坐标线性衰减，切向坐标随法向距离的平方根移动；
近端映射虽然不局部 Lipschitz，切向位移仍可求和。
在另一个凸模型中，同一二次目标的普通 PPA 线性收敛，
四次 Bregman 核却给出三次物理收敛，核的平坦性实际改变了隐式更新。
最后我们提出一个全局完成问题：局部吸引的紧光滑零簇能否成为全部零点？
对满足全尺度次线性反射界的完整普通近端映射，
局部指数与全局总指数的平衡给出欧拉限制；
弱凸径向目标中的极小点圆与额外临界点具体展示这一现象。

第~\ref{sec:rleb-gppa-main}~节建立物理空间的 GPPA 收敛定理和能量推论；
第~\ref{sec:deterministic-rleb-spine}~节给出普通 PPA 与固定正定度量的特例。
第~\ref{sec:examples}~节分析两个模型；
第~\ref{sec:topological-obstruction}~节讨论局部零簇的全域完成；
第~\ref{sec:comparison}~节识别线性阶条件并恢复经典收敛判据。
